(a) Since $\deg f>1$, all partial derivatives vanish at the origin. At a nonzero point of $X$, their simultaneous vanishing would make the corresponding point of $Y$ singular by \exref{I.5.8}. Thus the vertex is the only singular point.

(b) On the blowup chart with direction coordinate $t\ne0$, put $t=1$. The strict transform has equations $y=xu$, $z=xv$, and $f(1,u,v)=0$, so this chart is isomorphic to $\mathbb{A}^1\times(Y\cap\{t\ne0\})$. The latter is nonsingular: the gradient of $f(1,u,v)$ is nonzero at every point of its zero set. The other two charts are the same after permuting coordinates.

(c) Globally the strict transform consists of pairs $((x,y,z),(t:u:v))$ such that $f(t,u,v)=0$ and $(x,y,z)$ is proportional to $(t,u,v)$. Above the vertex the proportionality conditions impose no further restriction, so the exceptional fiber is $Z(f(t,u,v))\cong Y$.
